% =============================================================================
% APPENDIX AV: CORE-READY: The Willingness Function W(·)
% =============================================================================
% Module: BCM2_12_WIL (Willingness)
% Category: CORE
% Version: 1.0 (January 2026)
% Status: Final
% Language: English
%
% PURPOSE: This appendix answers the question "Is the individual willing and
% able to act?" It formalizes the Willingness Function W(·) that determines
% whether awareness translates into action.
%
% 10C POSITION: This is CORE 7 of 10 - the "READY" question.
%
% =============================================================================

\begin{tcolorbox}[colback=gray!5!white, colframe=gray!75!black,
    title=Appendix Header: AV CORE-READY]

\begin{tabular}{@{}ll@{}}
\textbf{Code:} & AV \\
\textbf{Category:} & CORE (Fundamental Theory) \\
\textbf{Full Name:} & CORE-READY: The Willingness Function \\
\textbf{Module:} & BCM2\_12\_WIL (Willingness) \\
\textbf{Version:} & 1.0 (January 2026) \\
\textbf{Status:} & Final \\
\end{tabular}

\vspace{0.5em}
\textbf{Purpose:} Answers the seventh fundamental question:
\textit{``Is the individual willing and able to act?''}
Formalizes the Willingness Function $W(\cdot)$ that converts awareness into action probability.

\end{tcolorbox}

\begin{tcolorbox}[colback=purple!5!white, colframe=purple!75!black,
    title=Chapter Linkage]

\textbf{Primary Chapter:} Chapter 12: Willingness and Action

\textbf{The Fundamental Question:}
\begin{quote}
\textit{``Given awareness of options and consequences, is the individual willing
to act? What barriers prevent action despite awareness and positive utility?''}
\end{quote}

\textbf{Relationship to Other COREs:}
\begin{itemize}[nosep]
\item Requires: $A > 0$ (CORE-AWARE must be satisfied first)
\item Modifies action probability: $P(\text{act}) = W(\cdot) | A > 0$
\item Context-dependent: $W = W(\Psi)$ via CORE-WHEN (V)
\end{itemize}

\end{tcolorbox}

\chapter{CORE-READY: The Willingness Function}
\label{app:core-ready}

\begin{abstract}
This appendix answers the seventh fundamental question:
\textit{``Is the individual willing and able to act?''}
Awareness alone does not guarantee action. The Willingness Function $W(\cdot)$
captures motivational barriers, implementation intentions, action costs, and
self-efficacy that determine whether potential actions become actual behavior.
\end{abstract}

\begin{tcolorbox}[colback=blue!5!white, colframe=blue!75!black,
    title=Quick Reference: CORE-READY]

\textbf{Core Question:} Is the individual willing and able to act?

\textbf{Key Formula:}
\begin{equation}
P(\text{act}) = W(\cdot) \cdot \mathbf{1}_{A > \tau_A} \quad \text{where } W \in [0, 1]
\label{eq:ready-action}
\end{equation}

\textbf{Willingness Components:}
\begin{itemize}[nosep]
\item $W_{\text{mot}}$: Motivation (desire to act)
\item $W_{\text{imp}}$: Implementation (concrete plans)
\item $W_{\text{eff}}$: Self-efficacy (belief in ability)
\item $W_{\text{cost}}$: Action cost (hassle factor, inverted)
\end{itemize}

\textbf{Cross-References:}
AU (CORE-AWARE), AW (CORE-STAGE), V (CORE-WHEN)

\end{tcolorbox}

\section{The Fundamental Question}
\label{sec:ready-fundamental}

\begin{tcolorbox}[colback=red!5!white, colframe=red!75!black,
    title=The Fundamental Question]
\textbf{Given awareness of options and their value, what determines whether
the individual actually takes action? What are the intention-action gaps?}
\end{tcolorbox}

The gap between intention and action is well-documented:
\begin{itemize}
\item People who want to save don't open retirement accounts
\item Patients who understand treatment benefits don't take medication
\item Citizens who value voting don't go to the polls
\end{itemize}

\section{Core Theory: The Willingness Function}
\label{sec:ready-theory}

\subsection{Definition}

\begin{definition}[Willingness Function]
The \textbf{Willingness Function} $W: \mathcal{S} \to [0, 1]$ maps states to
action probabilities conditional on awareness:
\begin{equation}
W(\cdot) = f(W_{\text{mot}}, W_{\text{imp}}, W_{\text{eff}}, W_{\text{cost}})
\end{equation}
\end{definition}

\subsection{Components of Willingness}

\begin{table}[htbp]
\centering
\caption{Willingness Components}
\label{tab:ready-components}
\renewcommand{\arraystretch}{1.3}
\begin{tabular}{@{}clp{5cm}@{}}
\toprule
\textbf{Component} & \textbf{Name} & \textbf{Definition} \\
\midrule
$W_{\text{mot}}$ & Motivation & Desire/drive to act on awareness \\
$W_{\text{imp}}$ & Implementation & Presence of concrete action plans \\
$W_{\text{eff}}$ & Self-Efficacy & Belief in own ability to execute \\
$W_{\text{cost}}$ & Action Cost & Hassle factors, friction (inverted) \\
\bottomrule
\end{tabular}
\end{table}

\section{Axiom System}
\label{sec:ready-axioms}

\begin{tcolorbox}[colback=gray!5!white, colframe=gray!75!black,
    title=Axiom WIL-1: Awareness Prerequisite]
\textbf{Statement:}
\begin{equation}
W > 0 \Rightarrow A > \tau_A \quad \text{(cannot be willing without awareness)}
\end{equation}

\textbf{Interpretation:} Willingness requires awareness. One cannot be willing
to do something they are unaware of.
\end{tcolorbox}

\begin{tcolorbox}[colback=gray!5!white, colframe=gray!75!black,
    title=Axiom WIL-2: Threshold Effects]
\textbf{Statement:}
\begin{equation}
P(\text{act}) = \begin{cases}
W(\cdot) & \text{if } W > \tau_W \\
0 & \text{otherwise}
\end{cases}
\end{equation}

\textbf{Interpretation:} Action requires willingness to exceed a threshold.
Below-threshold willingness produces no action.
\end{tcolorbox}

\begin{tcolorbox}[colback=gray!5!white, colframe=gray!75!black,
    title=Axiom WIL-3: Implementation Intention Effect]
\textbf{Statement:}
\begin{equation}
W_{\text{with plan}} > W_{\text{without plan}} \quad \text{(Gollwitzer effect)}
\end{equation}

\textbf{Interpretation:} Concrete implementation intentions ("when-where-how")
substantially increase willingness and action probability.
\end{tcolorbox}

\begin{tcolorbox}[colback=gray!5!white, colframe=gray!75!black,
    title=Axiom WIL-4: Friction Sensitivity]
\textbf{Statement:}
\begin{equation}
\frac{\partial W}{\partial \text{friction}} < 0 \quad \text{(small frictions have large effects)}
\end{equation}

\textbf{Interpretation:} Even small hassle factors disproportionately reduce willingness.
Removing friction increases action more than adding incentives.
\end{tcolorbox}

\section{The Intention-Action Gap}
\label{sec:ready-gap}

\begin{equation}
\text{Gap} = \text{Intention} - \text{Action} = A \cdot U_{\text{eff}} - W \cdot A \cdot U_{\text{eff}}
\end{equation}

The gap exists when $W < 1$ despite $A > 0$ and $U_{\text{eff}} > 0$.

\textbf{Common Gap Causes:}
\begin{enumerate}
\item Low self-efficacy ($W_{\text{eff}} < 1$)
\item Missing implementation plans ($W_{\text{imp}} < 1$)
\item High hassle costs ($W_{\text{cost}}$ low)
\item Present bias affecting motivation ($W_{\text{mot}}$ decays with delay)
\end{enumerate}

\section{Integration with Other COREs}

\begin{tcolorbox}[colback=yellow!5!white, colframe=yellow!75!black,
    title=Integration: CORE-READY $\leftrightarrow$ Other COREs]
\textbf{With CORE-AWARE (AU):}
Willingness requires awareness as prerequisite; they are sequential in the BCJ.

\textbf{With CORE-STAGE (AW):}
Different BCJ phases have different willingness requirements and barriers.

\textbf{With CORE-WHEN (V):}
Context affects willingness: deadlines increase urgency ($W_{\text{mot}}$),
friction reduces action cost sensitivity.

\textbf{With CORE-EIT (IE):}
Interventions can target willingness components directly (implementation intentions,
friction reduction, self-efficacy building).
\end{tcolorbox}

\section{Summary}

\begin{tcolorbox}[colback=gray!5!white, colframe=gray!75!black,
    title=Appendix AV Summary: CORE-READY]

\textbf{Core Contribution:}
Formalizes the Willingness Function $W(\cdot)$ that converts awareness into
action probability, explaining the intention-action gap.

\textbf{Key Formula:}
\begin{equation}
P(\text{act}) = W(\cdot) \cdot \mathbf{1}_{A > \tau_A}
\end{equation}

\textbf{Key Insights:}
\begin{itemize}[nosep]
\item Awareness is necessary but not sufficient for action
\item Willingness has four components: motivation, implementation, efficacy, cost
\item Implementation intentions substantially increase action probability
\item Small frictions have disproportionately large effects on willingness
\end{itemize}

\end{tcolorbox}

\vspace{1em}
\hrule
\vspace{0.5em}

\noindent\textit{Cross-references:}
Appendix AU (CORE-AWARE), AW (CORE-STAGE), V (CORE-WHEN), IE (CORE-EIT).

\noindent\textit{Master Bibliography:} \texttt{bcm\_master.bib}

\nocite{bcm_master}

% =============================================================================
% END OF APPENDIX AV
% =============================================================================
